\documentclass[11pt,a4paper]{article}
\usepackage[margin=1in]{geometry}
\usepackage{amsmath,amssymb,amsfonts}
\usepackage{mathtools}
\usepackage{lmodern}
\usepackage[T1]{fontenc}
\usepackage[utf8]{inputenc}
\usepackage{textcomp}
\usepackage{microtype}
\usepackage{parskip}
\usepackage{enumitem}
\usepackage{array}
\usepackage{booktabs}
\usepackage{longtable}
\usepackage{hyperref}
\hypersetup{hidelinks,
  pdftitle={Paper SCBU -- Substrate-Cosmology Boundary Unification},
  pdfauthor={Allen Proxmire}}
\setlength{\parindent}{0pt}
\setlength{\parskip}{6pt plus 2pt minus 1pt}

\title{\vspace{-1cm}\Large\bfseries Paper SCBU --- Substrate-Cosmology Boundary Unification}
\author{Allen Proxmire}
\date{May 2026}

\begin{document}
\maketitle

\section*{Preamble --- What This Paper Does NOT Claim}

\begin{enumerate}[noitemsep]
\item This paper does \textbf{not} claim that the MOND acceleration scale $a_0$ and the ED-SC canonical operating point $\xi_\mathrm{canonical}$ are numerically identical or equal at substrate level; they are scales of different physical dimensions (acceleration vs.\ length).
\item It does \textbf{not} claim a substrate-derivation of the numerical value of $\xi_\mathrm{canonical}$ from $H_0$ analogous to Paper\_029's derivation of $a_0$; that derivation is OPEN.
\item It does \textbf{not} claim that the unified substrate-cosmology boundary structure proposed here is the \textbf{only} possible cross-arc unification of MOND and ED-SC content.
\item It does \textbf{not} introduce new substrate primitives; only one paper-specific postulate (P-Substrate-Cosmology-Unified) is introduced.
\item It does \textbf{not} override or strengthen any constituent-paper verdict; it composes Paper\_028, Paper\_029, Paper\_096, and Paper\_097 into a single substrate-cosmology framing.
\item It does \textbf{not} claim novel empirical predictions; cross-arc structural identification is the deliverable, not new phenomenology.
\item ``Substrate-cosmology boundary'' = the substrate-level surface at $R_H = c/H_0$ (Paper\_028) beyond which cross-locus substrate influence cannot reach coherently. Both MOND and ED-SC content inherit projection structure from this boundary.
\end{enumerate}

\section*{Abstract}

The MOND acceleration scale $a_0 = cH_0/(2\pi)$ (Paper\_029) and the ED-SC canonical operating point $\xi_\mathrm{canonical} = 1.7575$ lu (Paper\_096) are anchored to substrate-cosmology boundary structure but have no current unifying paper in the corpus. This synthesis paper claims that both scales are projections of a single substrate-cosmology boundary at $R_H = c/H_0$ (Paper\_028): the MOND projection arises via azimuthal-Fourier on the residual $SO(2)$ symmetry of an accelerating chain (Paper\_029 \S5.1 mechanism); the ED-SC projection arises via cross-scale-invariance fixed-point matching on the kernel hierarchy (Paper\_096 mechanism). The two projections share the same substrate-cosmology origin and are FORCED to be structurally related under a single new postulate (P-Substrate-Cosmology-Unified), while remaining numerically distinct because they project the same boundary onto different observables (acceleration vs.\ substrate-length-unit). The synthesis prepares the corpus for ED-SC 4.x extensions to BH, MOND/galactic, Q-Compute, and soft-matter scales. The structural form is FORM-FORCED; the substrate-derivation of $\xi_\mathrm{canonical}$ from $H_0$ remains OPEN.

\section{Introduction}

\subsection{The cross-arc gap}

The ED corpus contains two independent canon-internal scale anchors of substrate-cosmological origin:

\begin{itemize}[noitemsep]
\item \textbf{Paper\_029} derives $a_0 = cH_0/(2\pi)$ as the substrate-level MOND transition acceleration. The factor $1/(2\pi)$ arises from azimuthal-Fourier-mode normalization on the residual $SO(2)$ symmetry of an accelerating chain projecting the cosmic decoupling surface (Paper\_028).
\item \textbf{Paper\_096} establishes the ED-SC cross-scale invariance with canon-internal canonical operating point $\xi_\mathrm{canonical} = 1.7575$ lu. This value is inherited from the ED-SC 3.3.x arc closure.
\end{itemize}

A cross-arc audit identified that no current paper unifies these two scales as projections of a common substrate-cosmology boundary structure. Both are anchored at $R_H = c/H_0$ in upstream content (Paper\_028), but the substrate-mechanism linking them across the MOND-arc and Wedges-arc has not been formalized. This paper supplies that formalization at verdict tier M3 per Paper\_095 methodology.

\subsection{Why unification matters}

The substrate-cosmology boundary $R_H = c/H_0$ is a load-bearing structural object in two arcs simultaneously:

\begin{itemize}[noitemsep]
\item \textbf{MOND arc:} $R_H$ is the cosmic decoupling surface (Paper\_028) supplying the geometric and dipole-projection structure that produces $a_0$.
\item \textbf{Wedges arc:} $R_H$ is the implicit far-IR endpoint of the three-regime RG flow (Paper\_097), beyond which cross-scale invariance ceases to apply at substrate level.
\end{itemize}

If both arcs invoke the same substrate-cosmology boundary, then their canon-internal scale anchors should be structurally related. The current corpus has the boundary referenced in both arcs but the \textbf{two anchors are not currently shown to inherit from it jointly}. This is a substrate-level structural gap, not an empirical gap.

Joint inheritance is a substrate-level structural claim, not a citation-hygiene issue --- it requires its own postulate and a derivation chain composing Paper\_028 + Paper\_029 + Paper\_096 + Paper\_097.

\subsection{How this paper fits in the corpus}

This paper sits at the intersection of the gravity-arc and wedges-arc, analogous to:

\begin{itemize}[noitemsep]
\item \textbf{Paper\_062} (cross-domain echo BH $\leftrightarrow$ Q-COMPUTE via shared V5 mechanism) --- methodological template.
\item \textbf{Paper\_071} (ER=EPR-class structural echo Entanglement $\leftrightarrow$ BH via shared V5 mechanism) --- methodological template.
\end{itemize}

Both Paper\_062 and Paper\_071 introduce single shared-mechanism postulates linking arcs that otherwise have independent derivation chains. This paper follows the same template: a single shared-mechanism postulate (P-Substrate-Cosmology-Unified) linking MOND arc and Wedges/RG arc via Paper\_028's substrate-cosmology boundary.

The structural-echo verdict in \S6 is \textbf{A$\to$position} --- same substrate boundary, two projection mechanisms --- not a new derivation of either $a_0$ or $\xi_\mathrm{canonical}$ from first principles.

\section{Primitive Inputs and Paper-Specific Postulates}

\subsection{Primitives invoked (per Paper\_087)}

\begin{itemize}[noitemsep]
\item \textbf{P03 (Channel + locus indexing; spatial homogeneity)} --- supplies substrate-level azimuthal and translation symmetry against which dipole-projection and cross-scale-invariance fixed-points are defined.
\item \textbf{P08 (Substrate scale $\ell_\mathrm{ED}$)} --- supplies the substrate length scale entering $\xi_\mathrm{canonical}$ via Paper\_096 inheritance.
\item \textbf{P11 (Commitment-irreversibility)} --- supplies the directional substrate evolution underlying the cosmic decoupling surface (Paper\_028) and the IR-regime cutoff (Paper\_097).
\item \textbf{P13 (Time homogeneity)} --- supplies $H_0$ as the cosmological rate parameter against which $R_H$ is defined.
\end{itemize}

\subsection{Upstream dependencies (I-rows)}

\begin{itemize}[noitemsep]
\item \textbf{I-028:} Cosmic decoupling surface $R_H = c/H_0$ (Paper\_028). \textbf{[LOAD-BEARING --- shared boundary source]}
\item \textbf{I-029:} $a_0 = cH_0/(2\pi)$ via dipole-projection on residual $SO(2)$ symmetry (Paper\_029). \textbf{[LOAD-BEARING --- MOND projection mechanism]}
\item \textbf{I-091:} Kernel cascade across scales (Paper\_091). \textbf{[LOAD-BEARING --- Wedges-arc scale-hierarchy framework]}
\item \textbf{I-096:} ED-SC cross-scale invariance with $\xi_\mathrm{canonical}$ canonical operating point (Paper\_096). \textbf{[LOAD-BEARING --- ED-SC projection mechanism]}
\item \textbf{I-097:} Three-regime RG with canon-internal 0.6 transition exponent (Paper\_097). \textbf{[LOAD-BEARING --- IR-regime endpoint]}
\item \textbf{I-037:} $a_0$ cosmological-rate invariance under local rescaling (Paper\_037 P-H0-Cosmological-Invariant). \textbf{[CONTEXT]}
\item \textbf{I-062:} Cross-domain echo via shared substrate mechanism (Paper\_062). \textbf{[CONTEXT --- methodological template]}
\item \textbf{I-071:} ER=EPR-class structural echo (Paper\_071). \textbf{[CONTEXT --- methodological template]}
\item \textbf{I-095:} Form-FORCED / value-INHERITED methodology with three-tier verdict classification (Paper\_095). \textbf{[METHODOLOGY]}
\end{itemize}

\subsection{Paper-specific postulates}

\textbf{P-Substrate-Cosmology-Unified:} The MOND acceleration scale $a_0 = cH_0/(2\pi)$ (Paper\_029) and the ED-SC canonical operating point $\xi_\mathrm{canonical}$ (Paper\_096) are two projections of a single substrate-cosmology boundary structure anchored at $R_H = c/H_0$ (Paper\_028). The MOND projection is azimuthal-Fourier on residual $SO(2)$ symmetry of an accelerating chain (Paper\_029 \S5.1 mechanism); the ED-SC projection is cross-scale-invariance fixed-point on the kernel hierarchy (Paper\_096 mechanism). The two projections share the same substrate-cosmology origin, are FORCED to be structurally related (via the shared boundary), and are constrained to scale jointly under $H_0$ variation. Numerical values of $a_0$ and $\xi_\mathrm{canonical}$ remain INHERITED --- the postulate fixes the structural relationship, not the numerical magnitudes.

This is the single new postulate introduced by this paper. It follows the cross-domain-echo template of Paper\_062 (P-V5-Shared-Mechanism) and Paper\_071 (P-V5-Shared-EntanglementGravity): a single load-bearing postulate that claims a shared substrate object underlies otherwise-independent arc derivations.

\section*{\S2.5 Load-Bearing Step Audit Table}

\begin{center}
\footnotesize
\begin{tabular}{cp{0.40\textwidth}p{0.18\textwidth}p{0.27\textwidth}}
\toprule
\textbf{\#} & \textbf{Step} & \textbf{Label} & \textbf{Notes} \\
\midrule
1 & Substrate-cosmology boundary $R_H = c/H_0$ & I & Paper\_028. \\
2 & $a_0 = cH_0/(2\pi)$ via dipole-projection & I & Paper\_029. \\
3 & ED-SC cross-scale invariance with $\xi_\mathrm{canonical}$ & I & Paper\_096. \\
4 & Three-regime RG with 0.6 transition exponent & I & Paper\_097. \\
5 & Memory-kernel cascade across scales & I & Paper\_091. \\
6 & Both $a_0$ and $\xi_\mathrm{canonical}$ anchored to $R_H$ & P-Substrate-Cosmology-Unified & Postulate. \\
7 & MOND projection: azimuthal-Fourier on residual $SO(2)$ & D-via-I & Paper\_029 \S5.1 mechanism. \\
8 & ED-SC projection: cross-scale-invariance fixed-point & D-via-I & Paper\_096 P-S1-Invariant mechanism. \\
9 & Two projections share substrate-cosmology origin & D & Composition of steps 6--8. \\
10 & Joint scaling under $H_0$ variation & D & Composition of step 9 + I-037 invariance. \\
11 & Numerical values of $a_0$ and $\xi_\mathrm{canonical}$ & I & Empirical via Paper\_029 + canon-internal via Paper\_096. \\
12 & Substrate-derivation of $\xi_\mathrm{canonical}$ from $H_0$ & OPEN & Not claimed. \\
13 & Cross-arc echo structural-positive verdict & A$\to$position & Composite. \\
\bottomrule
\end{tabular}
\end{center}

Verdict tier per Paper\_095: \textbf{M3 --- form-FORCED + value-INHERITED.} Potential upgrade to M2 (Intermediate Path C) if step 12 closes --- i.e., if a substrate-derivation of $\xi_\mathrm{canonical}$ from $H_0$ is supplied.

\section{The Unification}

\subsection{The shared substrate-cosmology boundary}

Paper\_028 establishes the cosmic decoupling surface at $R_H = c/H_0$ as the substrate-level boundary beyond which cross-locus substrate influence cannot reach a local chain coherently. This boundary is a structural property of the substrate participation graph under P11 commitment-irreversibility + P13 time homogeneity + the kernel-propagation rate $c$. The substrate-level identity:
\[
R_H = c/H_0
\]
is the canonical substrate-cosmology boundary. Two distinct projections of this boundary appear in the corpus.

\subsection{The MOND projection (Paper\_029 mechanism)}

By Paper\_029, an accelerating chain breaks the full $SO(3)$ rotational symmetry of substrate-graph adjacency to the residual $SO(2)$ symmetry about the acceleration axis. The substrate-graph response to the cosmic decoupling surface projects this content via azimuthal-Fourier-mode normalization. The leading anisotropic mode is the dipole ($l=1$), and the $|m|=1$ Fourier projection on the residual $SO(2)$ carries the geometric factor $1/(2\pi)$:
\[
a_0 = \frac{cH_0}{2\pi}.
\]
The factor $1/(2\pi)$ is FORM-FORCED by the azimuthal-Fourier normalization on the circle of period $2\pi$. The numerical $a_0 \approx 1.08 \times 10^{-10}\,\mathrm{m/s^2}$ at $H_0 \approx 70$ km/s/Mpc inherits Hubble-rate observational content.

\subsection{The ED-SC projection (Paper\_096 mechanism)}

By Paper\_096, the substrate kernel hierarchy (Paper\_092 three-index classification) admits a cross-scale invariance under scale transformations within the hydrodynamic-window-extended regime. The substrate-level S1 kernel-content is invariant under rescaling; the S2 ensemble-only content is regime-specific. The cross-scale-invariance fixed point is the canonical operating point:
\[
\xi_\mathrm{canonical} = 1.7575\,\mathrm{lu}.
\]
The numerical value is canon-internal --- inherited from ED-SC 3.3.x arc closure. At substrate level, $\xi_\mathrm{canonical}$ is anchored to the substrate-cosmology boundary $R_H = c/H_0$ via the IR-regime endpoint of the three-regime RG flow (Paper\_097): beyond $R_H$, cross-scale invariance ceases to apply at substrate level, and $\xi_\mathrm{canonical}$ is the canonical operating point at which the invariance is maximally saturated within the IR regime.

\subsection{Joint inheritance from $R_H$}

By P-Substrate-Cosmology-Unified, both $a_0$ and $\xi_\mathrm{canonical}$ inherit structurally from the same substrate-cosmology boundary $R_H = c/H_0$. The two projections are mechanistically distinct:

\begin{itemize}[noitemsep]
\item \textbf{MOND projection} projects $R_H$ onto \textbf{accelerating-chain content} via residual $SO(2)$ Fourier.
\item \textbf{ED-SC projection} projects $R_H$ onto \textbf{substrate-kernel-hierarchy content} via cross-scale-invariance fixed point.
\end{itemize}

But they share the same substrate-cosmology origin: $R_H$ is the load-bearing structural object in both cases. The structural relationship between the two projections is FORCED by joint inheritance from a single boundary.

\subsection{Joint scaling under $H_0$ variation}

A direct consequence of joint inheritance: under variation of $H_0$ (e.g., Hubble-tension resolution at any value in the empirical range), both $a_0$ and $\xi_\mathrm{canonical}$ must scale jointly. Specifically:

\begin{itemize}[noitemsep]
\item $a_0$ scales linearly with $H_0$: $a_0 = cH_0/(2\pi)$ --- direct.
\item $\xi_\mathrm{canonical}$ scales with $H_0$ through its anchoring to $R_H = c/H_0$ --- indirect.
\end{itemize}

The exact functional form of the $\xi_\mathrm{canonical}(H_0)$ scaling is not derived in this paper (OPEN at step 12 of the audit table); only the structural claim that the two scales co-vary under $H_0$ is FORCED. If empirical $H_0$ measurements pin down a value in the Hubble-tension range $[60, 80]$ km/s/Mpc, $a_0$ is predicted in the range $[0.93, 1.24] \times 10^{-10}$ m/s$^2$ and $\xi_\mathrm{canonical}$ should scale correspondingly. Independent variation of one without the other would refute P-Substrate-Cosmology-Unified.

\subsection{What this is NOT}

The unification is \textbf{structural}, not operational:

\begin{itemize}[noitemsep]
\item It does \textbf{not} claim that $a_0$ and $\xi_\mathrm{canonical}$ are numerically equal --- they have different physical dimensions.
\item It does \textbf{not} claim that the MOND mechanism and the ED-SC mechanism are operationally equivalent --- they project the same boundary via different operations (azimuthal-Fourier vs.\ cross-scale-invariance fixed-point).
\item It does \textbf{not} supply a first-principles substrate-derivation of $\xi_\mathrm{canonical}$ from $H_0$ --- only the inheritance pattern is forced; the substrate-derivation analogous to Paper\_029's $a_0$ derivation is OPEN.
\end{itemize}

This is \textbf{A$\to$position} at the composite-verdict level (audit-table row 13), following the cross-domain-echo template of Paper\_062 \S3.4 and Paper\_071 \S3.4.

\section{Distinguishing Structural vs.\ Numerical Content}

\textbf{Structural (D-via-I):}
\begin{itemize}[noitemsep]
\item Joint inheritance of $a_0$ and $\xi_\mathrm{canonical}$ from $R_H$.
\item Mechanism-distinctness of the two projections (azimuthal-Fourier vs.\ cross-scale-invariance fixed-point).
\item Joint scaling of $a_0$ and $\xi_\mathrm{canonical}$ under $H_0$ variation.
\end{itemize}

\textbf{Inherited (I):}
\begin{itemize}[noitemsep]
\item Numerical value of $a_0$ (via Paper\_029).
\item Numerical value of $\xi_\mathrm{canonical}$ (canon-internal via Paper\_096).
\item Numerical value of $H_0$ (empirical cosmological observation).
\item Numerical value of 0.6 transition exponent (canon-internal via Paper\_097).
\end{itemize}

\textbf{OPEN:}
\begin{itemize}[noitemsep]
\item Substrate-derivation of $\xi_\mathrm{canonical}$ from $H_0$ analogous to Paper\_029's derivation of $a_0$.
\item Quantitative functional form of $\xi_\mathrm{canonical}(H_0)$ scaling.
\item Extension to additional substrate-cosmology-anchored scales (BH horizon $r_H = 2GM/c^2$, Q-Compute $\mathcal{M}_\mathrm{crit}$, soft-matter $L_\mathrm{flow}$).
\end{itemize}

\section{Falsification Criteria}

\begin{itemize}[noitemsep]
\item \textbf{F1:} Empirical evidence that $a_0$ and $\xi_\mathrm{canonical}$ scale \textbf{independently} under $H_0$ variation --- refutes P-Substrate-Cosmology-Unified joint-inheritance claim. This is the sharpest falsifier: if Hubble-tension resolution moves $H_0$ to a specific value and only one of the two anchors tracks the change, the unified-source claim fails.
\item \textbf{F2:} Independent substrate-derivation of $\xi_\mathrm{canonical}$ that does NOT invoke $H_0$ or substrate-cosmology --- refutes the boundary-unification mechanism. If $\xi_\mathrm{canonical}$ is shown to derive from a purely substrate-internal mechanism (e.g., V1 kernel-width self-consistency at substrate scale alone), then the cross-arc unification dissolves.
\item \textbf{F3:} Refutation of Paper\_028 (cosmic decoupling surface), Paper\_029 ($a_0$ formula), or Paper\_096 (cross-scale invariance) propagates upward to refute this paper. The synthesis collapses if any single upstream pillar is refuted.
\item \textbf{F4:} Empirical or substrate evidence that a third or fourth canon-internal scale anchor (e.g., BH horizon $r_H$, Q-Compute $\mathcal{M}_\mathrm{crit}$, NS-Q operating point) is \textbf{not} anchored to $R_H = c/H_0$ --- would force re-examination of the substrate-cosmology-boundary framework's claimed cross-arc reach.
\item \textbf{F5:} Discovery of a fifth independent projection mechanism on $R_H$ producing a third distinct cosmologically-anchored scale not currently in the corpus --- would extend the framework rather than refute it.
\end{itemize}

\section{Position Statement}

The MOND acceleration scale $a_0 = cH_0/(2\pi)$ (Paper\_029) and the ED-SC canonical operating point $\xi_\mathrm{canonical} = 1.7575$ lu (Paper\_096) are \textbf{FORCED to share} a single substrate-cosmology boundary origin at $R_H = c/H_0$ (Paper\_028) under P-Substrate-Cosmology-Unified. The two projections are mechanistically distinct (azimuthal-Fourier on residual $SO(2)$ for MOND; cross-scale-invariance fixed-point on kernel hierarchy for ED-SC) but inherit structurally from the same boundary. Joint scaling under $H_0$ variation is FORCED; numerical values remain INHERITED.

The result is at verdict tier M3 per Paper\_095 methodology. Upgrade to M2 (Intermediate Path C) requires substrate-derivation of $\xi_\mathrm{canonical}$ from $H_0$ analogous to Paper\_029's derivation of $a_0$ --- OPEN.

This synthesis closes the cross-arc harmonization gap between the MOND arc and the Wedges arc. With this paper, \textbf{substrate-cosmology boundary structure is globally consistent across MOND arc, Wedges arc, and the cross-scale-invariance framework}. The result joins Paper\_062 (BH $\leftrightarrow$ Q-COMPUTE cross-domain echo) and Paper\_071 (Entanglement $\leftrightarrow$ BH ER=EPR-class echo) as the third corpus-level cross-arc synthesis paper at structural-echo verdict.

\end{document}
